

\begin{abstract}
The \textsc{Tensor Isomorphism} problem (\textsc{TI}) has recently emerged as having connections to multiple areas of research within complexity and beyond, but the current best upper bound is essentially the brute force algorithm. Being an algebraic problem, \textsc{TI} (or rather, proving that two tensors are \emph{non}-isomorphic) lends itself very naturally to algebraic and semi-algebraic proof systems, such as the Polynomial Calculus (PC) and Sum of Squares (SoS). For its combinatorial cousin \textsc{Graph Isomorphism}, essentially optimal lower bounds are known for approaches based on PC and SoS (Berkholz \& Grohe, SODA '17). Our main results are an $\Omega(n)$ lower bound on PC degree or SoS degree for \textsc{Tensor Isomorphism}, and a nontrivial upper bound for testing isomorphism of tensors of bounded rank.

We also show that PC cannot perform basic linear algebra in sub-linear degree, such as comparing the rank of two matrices, or deriving $BA=I$ from $AB=I$. As linear algebra is a key tool for understanding tensors, we introduce a strictly stronger proof system, PC+Inv, which allows as derivation rules all substitution instances of the implication $AB=I \rightarrow BA=I$. We conjecture that even PC+Inv cannot solve  \textsc{TI} in polynomial time either, but leave open getting lower bounds on PC+Inv for any system of equations, let alone those for \textsc{TI}. We also highlight many other open questions about proof complexity approaches to \textsc{TI}.
\end{abstract}



\section{Introduction}
Tensors have rapidly emerged as a fundamental data structure and key mathematical object of the 21st century. They play key roles in many different areas of science, engineering, and mathematics, from quantum mechanics and general relativity to neural networks \cite{TNN} and mechanical engineering. They arise in theoretical computer science in many ways, including from (post-quantum) cryptography \cite{Pat96, JQSY19}, derandomization, \textsc{Matrix Multiplication}, \textsc{Graph Isomorphism} \cite{GQ}, and several different parts of Geometric Complexity Theory. 

The fundamental notion of equivalence between tensors is that of isomorphism: two tensors are isomorphic if one can be transformed into the other by an invertible linear change of basis in each of the corresponding vector spaces. For example, two 2-tensors (=matrices) $M,M'$ are equivalent under this notion if there are invertible matrices $X,Y$ such that $XMY = M'$; similarly, two 3-tensors, represented by 3-way arrays $T_{ijk}, T'_{ijk}$ are isomorphic if there are three invertible matrices $X,Y,Z$ such that 
\begin{equation} \label{eq:TI}
\sum_{ijk} X_{ii'} Y_{jj'} Z_{kk'} T_{ijk} = T'_{i'j'k'}
\end{equation}
for all $i',j',k'$. The problem of (3-)\textsc{Tensor Isomorphism} (\textsc{TI}) is: given two such 3-way arrays, to decide if they are isomorphic. 

Over finite fields, two different versions of \textsc{TI} sandwich the complexity of its more famous cousin, \textsc{Graph Isomorphism}. Namely, as presented above, \textsc{GI} reduces to \textsc{TI}. In the other direction, over a finite field $\F=\F_{p^a}$, one can take an $n \times n \times n$ tensor and list it out ``verbosely'', as a set of $p^{an}$ many $n \times n$ matrices over $\F$; the isomorphism problem for such verbosely given tensors is equivalent to \textsc{Group Isomorphism} for a certain class of $p$-groups, widely believed to be the hardest cases of \textsc{Group Isomorphism} in general. As such, this verbose version of \textsc{TI} reduces to \textsc{GI}. Furthermore, with Babai's quasi-polynomial-time algorithm \cite{Babai}, the running times are quite close: $N^{O(\log N)}$ for \textsc{VerboseTI} and $N^{O(\log^2 N)}$ for \textsc{GI} (the exponent of the exponent was worked out by Helfgott \cite{helfgott}). Thus \textsc{TI} stands as a key obstacle to putting \textsc{GI} into $\mathsf{P}$.


In this paper, we initiate the study of (algebraic) proof complexity approaches to proving that two tensors are non-isomorphic.
Lower bounds on the Polynomial Calculus proof system imply lower bounds on Gr\"{o}bner basis techniques, and the latter are some of the leading methods for solving $\mathsf{TI}$-complete problems in cryptanalysis, e.g., \cite{DBLP:conf/eurocrypt/TangDJPQS22, DBLP:conf/eurocrypt/FaugereP06}.
In the context of \textsc{GI}, proof complexity plays an important role, through its connection with the Weisfeiler--Leman (WL) algorithm. Although this algorithm does not, on its own, solve \textsc{GI} in polynomial time \cite{CFI}, it is a key subroutine in many of the best algorithms for \textsc{GI}, both in theory \cite{Babai} and in practice (see \cite{mckay, MP}). And the picture that has emerged is that some proof systems for \textsc{GI} are known to be equivalent in power to WL \cite{AM, BerkholzG15}, and some lower bounds on proof systems are closely related to lower bounds for WL \cite{SSC, DBLP:conf/soda/ODonnellWWZ14, BerkholzG15}. Versions of WL for groups, and in particular finite $p$-groups---and hence, by the connection above, tensors over finite fields---have only recently begun to be explored \cite{BGLQW, BS1, BS2, CL}.

\subsection{Main results}

We focus on the Polynomial Calculus (PC, or Gr\"{o}bner) proof system \cite{CEI}, though our results will also hold for semi-algebraic proof systems such as Sum-of-Squares \cite{lasserre} as well. PC is used to show that a system of polynomial equations over a field $\F$ is unsatisfiable over the algebraic closure $\overline{\F}$, by deriving from the system of equations, in a line-by-line fashion, the contradiction $1=0$. The degree of a PC proof is the maximum degree of any line appearing in the proof, and it is a fundamental result that PC proofs of constant degree can be found in polynomial time \cite{CEI}. Much as WL informally ``captures all combinatorial approaches'' to \textsc{GI}, PC informally ``captures all approaches based on Gr\"{o}bner bases'' to showing that a system of polynomial equations is unsatisfiable. 

The systems of equations we study are, for two non-isomorphic tensors $T,T'$, the equations (\ref{eq:TI}) along with new matrices $X',Y',Z'$, and equations saying that these are the inverses of $X,Y,Z$, resp., viz.: $XX'=X'X = \Id$, and similarly for the others. The reason for introducing these new matrices, despite their not appearing in (\ref{eq:TI}), is that these invertibility equations are only degree 2. In contrast, if we instead used the determinant to indicate that $X$ was invertible, then our starting equations would have degree $n+1$, rather than constant degree $\leq 3$. Since the main complexity measure we study on PC is degree, having starting equations of degree $n$ would make it difficult to make meaningful lower bound statements.

Our first main result is (two proofs of) a lower bound on such techniques.

\begin{theorem} \label{thm:lowerMain}
Over any field, there are instances of $n \times n \times n$ \textsc{Tensor Isomorphism} that require PC degree $\Omega(n)$ to refute. Over $\R$, they also require Sum-of-Squares degree $\Omega(n)$ to refute.
\end{theorem}

The preceding goes by reduction from known lower bounds on PC for \textsc{Graph Isomorphism} \cite{BerkholzG15, BerkholzG16}, but has the disadvantage (from the tensor point of view) that the resulting tensors are quite sparse: in one direction, one of the slices is supported on an $\Omega(n) \times n$ matrix and all the others slices have support size 1. In a second proof (Section~6 (in the original article)), we get a polynomially worse lower bound $\Omega(\sqrt[4]{n})$, but with a reduction from \textsc{Random 3XOR} that is more direct. Indeed, we show that 3XOR itself can be viewed as a particular instance of a tensor problem \emph{without} gadgets; gadgets are only then needed to reduce from that tensor problem to \textsc{Tensor Isomorphism} itself. In contrast, the lower bounds on PC for \textsc{GI} (\emph{ibid.}) already use the Cai--F\"{u}rer--Immerman gadgets \cite{CFI} to reduce from XOR-SAT, and then even further gadgets are needed to reduce from \textsc{GI} to \textsc{TI}. 

Our technical contributions in the above theorem are thus three-fold: 
\begin{enumerate}
\item We show that the known reductions from \textsc{GI} to \textsc{TI} can be carried out in low-degree PC; 
\item We realize 3XOR very naturally as a tensor problem; and 
\item We give new reductions from 3XOR, through a series of tensor-related problems, to \textsc{TI}, that work as many-one reductions of the decision problems that can be carried out in low-degree PC.
\end{enumerate}

Complementing our lower bound, we also show that tensors of low rank are comparatively easy to test for (non)-isomorphism. Here, one of our upper bounds is in the weaker Nullstellensatz proof system (giving a stronger upper bound than only a PC upper bound). In the Nullstellensatz proof system, a proof that a system of equations $f_1 = \dotsb = f_m = 0$ is unsatisfiable consists of polynomials $g_i$ such that $\sum g_i f_i = 1$, and the Nullstellensatz degree is the maximum degree of any $g_i f_i$. The PC degree is always at most the Nullstellensatz degree, and the gap between the two can be nearly maximal for Boolean equations: $O(1)$ versus $\Omega(n / \log n)$ \cite{BCIP}. (For Boolean equations, there is always an $O(n)$ upper bound, though this does not apply to \textsc{TI}, see Remark~\ref{rmk:boolean} below).

\begin{theorem} \label{thm:upper}
Over any field, the Nullstellensatz degree of refuting isomorphism of two $n \times n \times n$ tensors of tensor rank $\leq r$ is at most $2^{O(r^2)}$. If working over a finite field $\F_q$ and including the equations $x^q - x$, the PC degree is at most $O(qr^2)$. 

In particular, isomorphism of constant-rank tensors can be decided in polynomial time.
\end{theorem}

\begin{remark} \label{rmk:boolean}
In many settings in proof complexity, Boolean axioms such as $x_i^2 = x_i$ or $x_i^2 = 1$ are included among the system of equations, and all such unsatisfiable systems of equations can be refuted in degree $O(n)$ ($n=$\# variables). If this were the case here, the above would only be interesting for very small values of $r$. In contrast, the equations for \textsc{TI} do not include any such Boolean axioms, and as such the naive degree upper bound is exponential in the number of variables. For $n \times n \times n$ tensors, this gives an upper bound of $2^{O(n^2)}$ \cite{sombra}, and thus, Theorem~\ref{thm:upper} gives nontrivial upper bounds all the way up to $r \leq n$. (We note that $n \times n \times n$ tensors can have rank up to $\Theta(n^2)$ \cite{lickteigTypical}.) The proof of Theorem~\ref{thm:upper} shows that for rank-$r$ tensors, \textsc{TI} can essentially be reduced to a system of equations in only $O(r^2)$ variables.
\end{remark}

\begin{remark} \label{rmk:rank}
For fixed $r$, testing if an $n \times n \times n$ tensor has rank $\leq r$ can be done in polynomial time, as follows. This will show that the algorithm of Theorem~\ref{thm:upper} genuinely solves the decision problem, and not just a promise problem. Given an $n \times n \times n$ tensor $T$, consider its three $n \times n^2$ flattenings. Use Gaussian elimination to put each such flattening, separately, into reduced row echelon form. If any of these flattenings has rank $> r$, reject. Otherwise, we get from this a list of $3r$ vectors $u_1, \dotsc, u_r, v_1, \dotsc, v_r, w_1, \dotsc, w_r$, such that $T$ lives in the $r \times r \times r$-dimensional space $\Span\{u_1, \dotsc, u_r\} \otimes \Span\{v_1, \dotsc, v_r\} \otimes \Span\{w_1, \dotsc, w_r\}$. Now in this space we can write down the Brent equations \cite{brent} for $T$ to have rank $\leq r$, which will be $r^3$ cubic equations in $3r^2$ variables (Brent's equations \cite[(5.06)]{brent} were specifically for the matrix multiplication tensor, but analogous equations are easily constructed for arbitrary tensors using the same idea). Since $r$ is constant, these equations may be solved in polynomial time (here we assume that we are either working over a finite field, a finite-degree extension of the rationals---see, for example, Grigoriev \cite{grigoriev}---or in the BSS model over an arbitrary field).
\end{remark}

Lastly, one may wonder why we focus on 3-\textsc{Tensor Isomorphism}, and not some of its many related variants. Indeed, just as there are other equivalence notions for matrices---such as conjugacy $XMX^{-1}$ and congruence $XMX^T$---there are many different kinds of multilinear objects that can be represented by multi-way arrays, including tensors, homogeneous polynomials (commutative or noncommutative), alternating matrix spaces, multilinear maps, and so on, each with their own corresponding notion of isomorphism. While these problems are indeed distinct, they are all equivalent under polynomial-time isomorphisms \cite{FGS, GQ}; such problems are called $\mathsf{TI}$-complete. Even isomorphism of $k$-way tensors (for any fixed $k \geq 3$) is equivalent to isomorphism of 3-tensors \cite{GQ}. This partially justifies our focus on \textsc{3-Tensor Isomorphism}. In the course of proving our reductions for the results stated above, we use many of the gadgets from \cite{FGS, GQ}, and show that such uses also often yield proof complexity reductions as well. Because of the variety of gadgets used in our reductions, we believe that many, if not all, of the gadgets from those results would also yield proof complexity reductions, so the proof complexity of all the known $\mathsf{TI}$-complete problems should be polynomially related.

\subsection{Comparison with linear algebra, a new proof system, and a conjecture} \label{sec:conj}
As linear algebra is part of the core toolkit for understanding tensors, it is natural to wonder how linear algebra can help in algebraic proof complexity approaches to \textsc{TI}. We believe that even if it had the ``full power'' of linear algebra at its disposal ``for free,'' PC could still not solve \textsc{TI} efficiently. We begin to make this precise in this section.

Some basic derivations in linear algebra are to relate the ranks of two matrices and to derive $BA=I$ from $AB=I$ (the Inversion Principle, one of the so-called ``hard matrix identities'' \cite{SC}, only recently shown to have short $\mathsf{NC}^2$-Frege proofs \cite{HT}). Soltys \cite{soltys} and Soltys \& Cook \cite{SC} discuss the relationship between these and other standard implications in linear algebra. We show that PC is not strong enough to prove these in low-degree:

\begin{theorem}
The unsatisfiable system of equations $XY=\Id_n$ where $X$ is $n \times r$ and $Y$ is $r \times n$ with $1 \leq r < n$, requires degree $\geq r/2+1$ to refute in PC, over any field.
\end{theorem}
We refer to this system of equations as the Rank Principle, as refuting them amounts to showing that $\rk \Id_n > r$. 

\begin{theorem}
Any PC derivation of $BA = I$ from $AB=I$, where $A,B$ are $n \times n$ matrices with $\{0,1\}$ entries, requires degree $\geq n/2+1$, over any field.
\end{theorem}

We also observe that the Rank Principle can be derived in low degree from the Inversion Principle. 

Although it remains open whether the Inversion Principle is ``complete'' for linear-algebraic reasoning (see \cite{soltys, SC}), we introduce the proof system PC+Inv in an attempt to capture some linear-algebraic reasoning that seems potentially useful for \textsc{TI}. PC+Inv has all the same derivation rules as PC, but in addition, for any square matrices $A,B$ (whose entries may themselves be polynomials---that is, we allow substitution instances), we have the rule
\[
\frac{AB=I}{BA=I}.
\]
where the antecedent represents the set of $n^2$ equations corresponding to $AB=I$, and similarly the consequent denotes the set of $n^2$ equations $BA=I$ (see \ref{sec:prelim:inv} for more details). Degree is still measured in the usual way, but this rule lets us ``cut out'' the high-degree proof that would usually be required to derive $BA=I$ from $AB=I$. We now formalize our intuition that linear algebra should not suffice to solve \textsc{TI} efficiently in the following:

\begin{conjecture} \label{conj:inv}
\textsc{Tensor Isomorphism} for $n \times n \times n$ tensors requires degree $\Omega(n)$ in PC+Inv, over any field.
\end{conjecture}

Despite the conjecture, we do not yet know how to prove lower bounds on PC+Inv for any unsatisfiable system of equations, let alone those coming from \textsc{TI}. Mod $p$ counting principles (for $p$ different from the characteristic of the field) strike us as potentially interesting instances to examine for PC+Inv lower bounds, before tackling a harder problem like \textsc{TI}. In the final section, we highlight many other open questions around the proof complexity of \textsc{TI}.

\subsection{Organization}
In Section~\ref{sec:prelim} we cover preliminaries. In Section~3 (in the original article) we prove the lower bounds on linear algebraic principles just discussed. In Section~4 (in the original article) we prove the upper bound for isomorphism of bounded rank tensors (Theorem~\ref{thm:upper}). In Section~\ref{sec:lowerboundGI} we prove Theorem~\ref{thm:lowerMain} by reduction from \textsc{GI}. In Section~6 (in the original article) we prove the polynomially related lower bound by direct reduction from \textsc{Random 3XOR}.
